\subsection{QEMU Object Model (QOM)}
\label{chap:QEMU_QOM}

The QEMU Object Model (QOM) subsystem is QEMU's object-oriented framework. 
This subsystem provides a device modeling abstraction layer by implements a 
type system using pure C which, despite C's lack of native object-oriented programming 
(OOP) features, allows dynamic type registration, inheritance, and runtime 
polymorphism, providing\cite{Bellard2005Qemu} \cite{Qemu_QOM}:

\begin{itemize}
    \item \textbf{Object-oriented programming:    }Without requiring an OOP 
                                                    language like C++.

    \item \textbf{Dynamic type registration:      }Types can be registered at runtime without 
                                                    recompilation.
    
    \item \textbf{Inheritance and polymorphism:   }Types inherit from parent types (single 
                                                    inheritance) and override methods.
    
    \item \textbf{Runtime reflection:             }All objects and properties are 
                                                    discoverable via standardized queries.
    
    \item \textbf{Dynamic properties:             }Each object instance has configurable, 
                                                    typed properties.
    
    \item \textbf{Composition tree:               }Hierarchical object relationships enabling complex 
                                                    system modeling.
\end{itemize}

Before the introduction of QOM, QEMU development was fragmented. Each 
board implemented its devices independently, requiring duplicate code; device 
configurations were hardcoded; and extending QEMU with new hardware was cumbersome. 
The introduction of the QOM subsystem resolved these problems by unifying 
device management through the implementation of:

\begin{itemize}
    \item \textbf{Uniform object model:             }All devices are Objects with properties.
    
    \item \textbf{Runtime type registration:        }Types defined via TypeInfo at startup.
    
    \item \textbf{Common instantiation interface:   }Devices created uniformly by type name.
    
    \item \textbf{Dynamic configuration:            }Properties accessible and modifiable via a 
                                                        standardized API.
\end{itemize}


QOM's type system is foundational to QEMU device modeling. Every 
emulated peripheral (timers, ADCs, I2C controllers) is a QOM Object instance with 
an ObjectClass defining its behavior. This abstraction decouples device logic from 
QEMU core, enabling researchers to develop and validate new device models—such 
as the STM32F405 I2C controller developed in this thesis—within a consistent, 
well-defined framework. The figure **F1** illustrates the scope of QOM's management 
within QEMU.

- **F1**

\subsubsection{QOM API Core Elements}

\textbf{ObjectClass:}\cite{Qemu_QOM_API} Analogous to a class in any OOP language, the C \\ 
\texttt{ObjectClass} structure represents the metadata that describe a type and that is 
shared by all instances of that type. It holds:

\begin{itemize}
    \item \textbf{Type information}: Unique identifier for runtime type checking
    \item \textbf{Virtual method pointers}: Functions that can be overridden by derived types
    \item \textbf{Interface implementations}: Lists of interfaces this type provides.
    \item \textbf{Shared class properties}: Default values and metadata for all instances
\end{itemize}

\textbf{Object:}\cite{Qemu_QOM_API} Every object in QEMU inherits from the C \texttt{Object} structure, which 
is the fundamental component of the entire type system. Each individual object 
is an \textbf{object instance} with its own state, which can be instantiated, configured, 
introspected, and destroyed. Each \texttt{Object} instance encapsulates:

\begin{itemize}
    \item \textbf{Type identification} via a class pointer
    \item \textbf{Reference counting} for memory management
    \item \textbf{Dynamic properties} storage and access
    \item \textbf{Composition tree relationships} (parent pointers) enabling hierarchical 
                                        organization
\end{itemize}

\textbf{TypeInfo:}\cite{Qemu_QOM_API} The C \texttt{TypeInfo} structure is the complete blueprint for registering 
a new QOM type at QEMU startup. This structure tells QEMU exactly how to create, initialize, 
and manage instances of the QOM type.

\begin{itemize}
    \item Type name and parent class (e.g., parent=\texttt{"sysbus-device"})
    \item Instance and class sizes (memory layout)
    \item Constructor (\texttt{instance\_init}), post-constructor (\texttt{instance\_post\_init}), and 
            destructor (\texttt{instance\_finalize}) callbacks
    \item Class initialization callback (\texttt{class\_init}) where virtual methods are assigned
\end{itemize}

This will be referenced later in XXXX, but it serves to exemplify the use of 
the \texttt{TypeInfo} structure. The \texttt{STM32I2CClass} is registered via \texttt{TypeInfo} with 
parent=\texttt{'sysbus-device'}, \texttt{instance\_size=sizeof(STM32I2CState)}, and 
\texttt{class\_init} \texttt{=stm32\_i2c\_class\_init}. This registration enables QEMU to instantiate the I$^2$C 
controller and configure its base address, interrupt lines, and master/slave modes at 
runtime.

\textbf{ObjectProperty:}\cite{Qemu_QOM_API} Each object's configurable properties—e.g., base addresses, 
enable bits, interrupt vectors—are registered via \texttt{ObjectProperty} descriptors. These 
define getters/setters, type metadata, and default values, enabling firmware 
configuration without recompilation.
